\documentclass[11pt]{article}
\usepackage[T1]{fontenc}
\usepackage[a4paper,margin=1in]{geometry}
\usepackage{amsmath,amssymb,amsthm}
\usepackage{mathrsfs}
\usepackage[hidelinks]{hyperref}
\usepackage{microtype}

\newtheorem{theorem}{Theorem}
\newtheorem{lemma}[theorem]{Lemma}
\newtheorem{corollary}[theorem]{Corollary}
\newtheorem{remark}[theorem]{Remark}

\newcommand{\HH}{\mathscr H^2}
\newcommand{\R}{\mathbb R}
\newcommand{\N}{\mathbb N}

\setlength{\parindent}{0pt}
\setlength{\parskip}{0.42em}

\title{A Smooth Recurrent-Spike Obstruction for\
Weighted Boundary Embeddings of Dirichlet Series}
\author{Run \texttt{fa\_banach\_001} --- lane 1}
\date{13 August 2026}

\begin{document}
\maketitle

\begin{center}
\fbox{\parbox{0.9\textwidth}{
\textbf{Status: candidate partial result.}
The source asks which nonnegative weights $\omega$ yield a bounded
critical-line embedding for the Hardy--Hilbert space of Dirichlet series.
We prove that the necessary condition $\omega\in L^1\cap L^\infty$ is not
sufficient, even for smooth weights tending to zero.  In fact, a bad weight
can have an admissible symmetric decreasing rearrangement.
}}
\end{center}

\section{Source problem and result}

Let $\HH$ be the Hilbert space of Dirichlet series
\[
 f(s)=\sum_{n\geq1}a_n n^{-s},\qquad
 \|f\|_{\HH}^2=\sum_{n\geq1}|a_n|^2.
\]
Brevig and Perfekt ask which nonnegative measurable $\omega$ satisfy
\begin{equation}\label{eq:embedding}
 \int_{\R}|f(1/2+it)|^2\omega(t)\,dt
 \leq C_\omega\|f\|_{\HH}^2
 \qquad(f\in\HH).
\end{equation}
They note that $L^1(\R)\cap L^\infty(\R)$ is necessary and prove
\eqref{eq:embedding} when $\omega(t)=v(|t|)$ with $v$ decreasing and
$v\in L^1\cap L^\infty$; see Lemma~7.2 and the ensuing open problem in
arXiv:2005.05094~\cite{source}.

\begin{theorem}\label{thm:main}
There is a nonnegative function $\omega\in C^\infty(\R)$ such that
\[
 \lim_{|t|\to\infty}\omega(t)=0,
 \qquad
 \omega\in L^q(\R)\quad(1\leq q\leq\infty),
\]
but \eqref{eq:embedding} fails.  Moreover, the symmetric decreasing
rearrangement $\omega^\#$ of $\omega$ satisfies \eqref{eq:embedding}.
Consequently admissibility is not rearrangement invariant; in particular,
it cannot be characterized by the distribution function of $\omega$ or by
rearrangement-invariant norms alone.
\end{theorem}

\section{Proof intuition}

The normalized truncated zeta kernel
\[
 F_N(s)=\frac1{\sqrt{H_N}}
 \sum_{n\leq N}\frac1{\sqrt n}\,n^{-s},
 \qquad H_N=\sum_{n\leq N}\frac1n,
\]
has $\HH$ norm one.  On the critical line it has height about
$\sqrt{\log N}$ on a window of width about $1/\log N$ around $t=0$.
The logarithms of the primes are rationally independent, so Kronecker's
theorem makes all phases $n^{-it}$, $n\leq N$, return simultaneously near
$1$ at arbitrarily many widely separated times.  Hence the same narrow,
high peak recurs as often as desired.

Place smooth bumps of height $h_k$ on $K_k$ such recurrences for a rapidly
growing cutoff $N_k$.  Their total $L^1$ cost is of order
$h_kK_k/\log N_k$, while testing against $F_{N_k}$ yields a contribution of
order $h_kK_k$.  Choosing $h_k=1/k$, $K_k=k^2$, and $N_k$ sufficiently
large makes the first quantities summable and the second tend to infinity.

\section{The recurrent peak}

\begin{lemma}\label{lem:peak}
Fix $N\geq2$.  There are arbitrarily large positive numbers $\tau$ such
that, whenever $|u|\leq(16\log N)^{-1}$,
\[
 |F_N(1/2+i(\tau+u))|^2\geq \frac{49}{64}H_N.
\]
Consequently one may choose any prescribed finite number of such $\tau$,
with the corresponding intervals pairwise disjoint and lying beyond any
previously fixed compact set.
\end{lemma}

\begin{proof}
Let $\mathcal P_N$ be the finite set of primes at most $N$, and put
$r_N=\lfloor\log_2N\rfloor$.  Unique factorization shows that
$\{\log p:p\in\mathcal P_N\}$ is linearly independent over $\mathbb Q$.
Kronecker's theorem therefore gives arbitrarily large $\tau>0$ for which
\[
 |p^{-i\tau}-1|\leq\frac1{16r_N}
 \qquad(p\in\mathcal P_N).
\]
Writing $n\leq N$ as a product of at most $r_N$ primes, counted with
multiplicity, and using
$|z_1\cdots z_m-1|\leq\sum_j|z_j-1|$ for unimodular $z_j$, gives
$|n^{-i\tau}-1|\leq1/16$.  If
$|u|\leq(16\log N)^{-1}$, then
\[
 |n^{-iu}-1|\leq |u|\log n\leq\frac1{16},
\]
and hence $|n^{-i(\tau+u)}-1|\leq1/8$ for every $n\leq N$.  It follows that
\begin{align*}
 \left|\sum_{n\leq N}n^{-1-i(\tau+u)}\right|
 &\geq H_N-
 \sum_{n\leq N}\frac{|n^{-i(\tau+u)}-1|}{n}\\
 &\geq\frac78H_N.
\end{align*}
Division by $\sqrt{H_N}$ proves the estimate.  The positive tail of the
dense Kronecker orbit is dense as well, so the returns can be chosen
successively beyond arbitrary compact sets.
\end{proof}

\section{Construction of the weight}

Choose integers $N_k\geq2$ so rapidly that
\begin{equation}\label{eq:Nchoice}
 \log N_k\geq 2^k k^2,
\end{equation}
and set
\[
 h_k=\frac1k,\qquad K_k=k^2,
 \qquad \delta_k=\frac1{16\log N_k}.
\]
Using Lemma~\ref{lem:peak}, choose $K_k$ recurrence centers
$\tau_{k,1},\ldots,\tau_{k,K_k}$ so that every interval
$(\tau_{k,j}-\delta_k,\tau_{k,j}+\delta_k)$ is disjoint from all the
others, and arrange the blocks successively toward $+\infty$.

Let $\eta\in C_c^\infty((-1,1))$ satisfy $0\leq\eta\leq1$ and
$\eta=1$ on $[-1/2,1/2]$, and define
\begin{equation}\label{eq:omega}
 \omega(t)=\sum_{k\geq1}h_k\sum_{j=1}^{K_k}
 \eta\left(\frac{t-\tau_{k,j}}{\delta_k}\right).
\end{equation}
The supports are disjoint and locally finite.  Thus $\omega$ is smooth,
$0\leq\omega\leq1$, and $\omega(t)\to0$ as $|t|\to\infty$ because
$h_k\to0$.

For every $1\leq q<\infty$, disjointness, a change of variables, and
\eqref{eq:Nchoice} give
\[
 \int_{\R}\omega(t)^q\,dt
 =\|\eta\|_{L^q}^q
 \sum_{k\geq1}K_kh_k^q\delta_k
 \leq \frac{\|\eta\|_{L^q}^q}{16}
 \sum_{k\geq1}\frac{k^{-q}}{2^k}<\infty.
\]
Hence $\omega\in L^q$ for all stated $q$, including $q=\infty$.

On the core interval
$|t-\tau_{k,j}|\leq\delta_k/2$, the $j$th bump equals $h_k$, and
Lemma~\ref{lem:peak} applies.  Since $H_N\geq\log(N+1)>\log N$,
\begin{align*}
 \int_{\R}|F_{N_k}(1/2+it)|^2\omega(t)\,dt
 &\geq K_k h_k\delta_k\frac{49}{64}H_{N_k}\\
 &\geq \frac{49}{1024}K_kh_k
 =\frac{49}{1024}k\longrightarrow\infty.
\end{align*}
But $\|F_{N_k}\|_{\HH}=1$, so \eqref{eq:embedding} fails.

Finally, $\omega^\#$ is an even radially decreasing nonnegative function
equimeasurable with $\omega$.  In particular it lies in
$L^1\cap L^\infty$.  Lemma~7.2 of the source paper applies to
$\omega^\#$ and proves its admissibility, completing the proof of
Theorem~\ref{thm:main}.

\section{A useful exact formulation and the remaining gap}

For $\omega\in L^1(\R)$, write
$\widehat\omega(\xi)=\int_{\R}\omega(t)e^{-it\xi}\,dt$.  Expanding a
Dirichlet polynomial shows that \eqref{eq:embedding} is equivalent to
uniform boundedness on $\ell^2$ of the finite sections of the positive
matrix
\begin{equation}\label{eq:gram}
 G_\omega(m,n)=
 \frac{\widehat\omega(\log(m/n))}{\sqrt{mn}}.
\end{equation}
This is an exact coefficient criterion, but it is not yet a geometric
description of the weights.

The source's local embedding theorem also immediately gives the broader
sufficient condition
\[
 \sum_{j\in\mathbb Z}
 \operatorname*{ess\,sup}_{t\in[j,j+1]}\omega(t)<\infty.
\]
The construction above shows why $L^1\cap L^\infty$ alone misses essential
information: arithmetic placement of thin pieces can synchronize with
simultaneous almost-periodic recurrences.  A full intrinsic
characterization remains open; the present theorem rules out the most
natural rearrangement-invariant answer.

\begin{thebibliography}{9}
\small
\bibitem{source}
O. F. Brevig and K.-M. Perfekt,
\emph{A mean counting function for Dirichlet series and compact composition
operators}, Adv. Math. 385 (2021), 107775; arXiv:2005.05094.

\bibitem{HLS}
H. Hedenmalm, P. Lindqvist, and K. Seip,
\emph{A Hilbert space of Dirichlet series and systems of dilated functions
in $L^2(0,1)$}, Duke Math. J. 86 (1997), 1--37.

\bibitem{OS}
J.-F. Olsen and E. Saksman,
\emph{On the boundary behaviour of the Hardy spaces of Dirichlet series and
a frame bound estimate}, J. Reine Angew. Math. 663 (2012), 33--66;
arXiv:0907.2708.
\end{thebibliography}

\end{document}
